A hypernetwork can condition a frozen language model by generating low-rank weight updates from sensor features. The same bridge also adds trainable capacity, so improved text likelihood alone is an ambiguous indicator of sensor use. We separate these effects with independently sampled text, paired category probes, and a separately trained instruction-generation cohort. The study retains the original 33.6M host and adds pinned 230M and 0.8B hosts, with matched constant, shuffled, random, static-adapter, and frozen references. Under independent sampling, 9 of 9 true-minus-constant BPB intervals include zero. This does not establish equivalence or extend to all paired likelihood objectives. Full held-out populations, grouped intervals, and same-checkpoint interventions assess conditioning. Numerical training failures affected 25 of the 411 declared trials; the corresponding endpoint estimates are unavailable. Mini reaches macro rank-1 accuracy 0.461 on audio, 0.582 on vision; Qwen 0.8B reaches macro rank-1 accuracy 0.838 on IMU, 0.421 on audio, 0.451 on vision. On held-out templates, Qwen's IMU joint label and sentence success is 0.838 and 0.689; sensor replacement tests their input dependence. Success on both reversed rule mappings is 0.005. Audio generation varies substantially across seeds. We report training-budget trajectories, effective-update variation, three composition operators, and measured bridge costs. The dense output head grows with adapted dimensions, and workstation results do not establish edge efficiency. Audio evidence is limited by the original encoder's use of held-out inputs in normalization. Claims about likelihood, input prefixes, and generation are scoped to the tested objectives and interfaces; the complete evidence and per-paper reproduction code accompany the revision.
